\section{A complete finite decision at each occurrence count}
\label{sec:decision}

The parameter bounds become useful only when accompanied by an exhaustive
and economical way to generate the relevant markings. Scanning all integers
below the height bound would obscure the main gain: the height is
exponential in $h$, but the number of Markov occurrences below it is
polynomial in $h$. We first establish this distinction, then explain how
to recognize the geometric meaning of every equality.

\subsection{Generating all minimum seeds}

\begin{lemma}[Minimum seeds and ordered occurrences]\label{lem:seed-completeness}
Every minimum-centered positive integral Markov ray with fixed parameter
$p$ has one of the two initial pairs
\begin{equation}\label{eq:two-seeds}
 (u_0,u_1)=(a,b)\quad\hbox{or}\quad(a,3pa-b),
\end{equation}
where $(a,b,p)$ is an ordered positive Markov triple. Conversely,
each such triple supplies both rays in \eqref{eq:two-seeds}, with
$u_0$ a minimum. All ordered occurrences with maximum at most $H$ are
generated by the two initial occurrences $(1,1,1),(1,1,2)$ and the tree
rooted at $(1,2,5)$ with children
\begin{equation}\label{eq:ordered-tree}
 (a,b,c)\longmapsto(a,c,3ac-b),\qquad(b,c,3bc-a),
\end{equation}
pruning a child when its maximum exceeds $H$.
\end{lemma}

\begin{proof}
Let $x=u_0$, $z=\min(u_{-1},u_1)$, and let $y$ be the other neighbor.
They are roots of $f(T)=T^2-3pxT+p^2+x^2$ and satisfy $x\leqslant z$.
The minimum bound implies $x\leqslant p$, while
$$
 f(p)=x^2-(3x-2)p^2\leqslant0.
$$
Thus $z\leqslant p\leqslant y$, and $(x,z,p)$ is an ordered Markov
triple. Conversely, for an ordered triple $(a,b,p)$ the other neighbor is
$$
 3pa-b=\frac{p^2+a^2}{b}\geqslant p\geqslant a.
$$
Both immediate neighbors are at least $a$; the recurrence, or
Lemma~\ref{lem:ray-model}, then makes $a$ a global minimum.
The two orientations are distinct: equality of the neighbors would give
$b^2=p^2+a^2>p^2$, contradicting $b\leqslant p$.

For completeness of the tree, consider an ordered triple $a\leqslant b\leqslant c$.
The conjugate root $c'=3ab-c=(a^2+b^2)/c$ is positive. The polynomial
$g(T)=T^2-3abT+a^2+b^2$ satisfies
$$
 g(b)=a^2-(3a-2)b^2\leqslant0,
$$
with equality only if $a=b=1$. Outside $(1,1,1)$ and $(1,1,2)$,
we have $0<c'<b<c$. Replacing the largest entry and reordering therefore
gives a uniquely determined parent of smaller height. Iteration reaches
the initial occurrences. Equal coordinates can occur only there:
descent preserves the multiset of pairwise greatest common divisors,
which are all one at the initial triples. Thus $a=b$ forces $a=b=1$,
while $g(b)<0$ excludes $b=c$ outside the initial cases.
Reversing this unique descent
gives exactly \eqref{eq:ordered-tree}, with strictly increasing maximum.
No uniqueness assertion concerning triples with the same maximum is used.
\end{proof}

The construction has a literal matrix realization. Set
$C=EFE=\left(\begin{smallmatrix}5&2\\2&1\end{smallmatrix}\right)$.
For the three initial occurrences use, respectively,
$$
 (U,V)=(E,F),\qquad(E,FE),\qquad(E,C).
$$
The normalized traces of $(U,V,UV)$ are the displayed triples.
The two child operations are
\begin{equation}\label{eq:basis-tree}
 (U,V)\longmapsto(U,UV),\qquad(V,UV).
\end{equation}
At a retained occurrence $(a,b,p)$, choose
$(A_0,B)=(U,(UV)^{-1})$ or $(U,UV)$. These realize exactly
\eqref{eq:two-seeds}. Starting from the root columns
$[E]=(1,0)$, $[F]=(0,1)$, update the columns by addition in
\eqref{eq:basis-tree}, and negate the second one when inverting $UV$.
Thus the representation and its marking are generated together.

This model loses no marked realization. Diagonalizing its second
generator determines the diagonal entries of the first from the complete
trace triple; its off-diagonal product is positive and fixed by the
determinant. The real conjugation in Lemma~\ref{lem:ray-model} therefore
identifies any two such full characters. It conjugates the whole group,
so also transports its isometries. Agreement of just one matrix trace
would not provide this conclusion.

\subsection{Counting occurrences rather than integer labels}

\begin{lemma}[A quadratic count in word length]\label{lem:seed-count}
If every retained positive Cohn word has length at most $J$, the seed
construction has at most $J(J-1)+4$ oriented seeds. A word of length
$j$ in $E,C$ has trace strictly larger than $(5/2)^j$.
\end{lemma}

\begin{proof}
The positive tree in \eqref{eq:basis-tree}, starting at $(E,C)$,
assigns each new maximum a positive primitive count vector $(r,s)$.
These vectors are distinct before traces are evaluated. Indeed adjacent
vectors have determinant of absolute value one, and inserting their sum
splits the interval between their slopes into two disjoint open intervals.
A later insertion stays strictly inside its own interval. Induction gives
the usual Farey mediant tree, without repeated interior slopes. At length
$j=r+s$ there are at most $j-1$ vectors. Summing and including the two
initial occurrences gives at most $2+J(J-1)/2$ occurrences; each has two
orientations.

Entrywise, $C\geqslant E>0$. Replacing every $C$ by $E$ can only decrease
a positive product and its trace. The expanding eigenvalue of $E$ is
$(3+\sqrt5)/2>5/2$, so $\tr(E^j)>(5/2)^j$.
\end{proof}

\begin{theorem}[Explicit complete decision domain]\label{thm:finite-decision}
Fix $h\geqslant2$. All equalities
$\tr A=\tr W_{h,k,t}(A,B)$, over every marked modular-torus basis,
$k\geqslant1$, and $1\leqslant t<h$ coprime to $h$, can be classified
by a finite algorithm. Define
\begin{equation}\label{eq:decision-heights}
 \begin{gathered}
 H_+=4^{h+1}-1,\qquad J_+=2h+3,\qquad S_+=J_+(J_+-1)+4,\\
 H_0=(100h^2)^{h+1}-1,\qquad
 L=\lceil\log_2(100h^2)\rceil,\\
 J_0=(h+1)L+1,\qquad S_0=J_0(J_0-1)+4.
 \end{gathered}
\end{equation}
At most
\begin{equation}\label{eq:decision-count}
 \begin{split}
 N_+(h)&=2S_+(3h-1)(h-1)(2h+1),\\
 N_0(h)&=2S_0(h-1)(3h+1)
 \end{split}
\end{equation}
exact trace comparisons suffice. Their sum is $O(h^5)$, using integers
with $O(h^3)$ bits. Every equal pair can also be tested for actual
isometry using twelve fixed integral homology actions.
The complexity is polynomial in the numerical value of $h$.
\end{theorem}

\begin{proof}[Construction of the domain]
Apply the centering and reversal of Section~\ref{sec:gap}.
The positive-product branch has $1\leqslant l\leqslant3h-1$ and
$p\leqslant H_+$ by Theorem~\ref{thm:positive-bound}. The minimum-product
branch has $l=0$ and $p\leqslant H_0$ by
Theorem~\ref{thm:minimum-bound}. Generate every seed up to the relevant
height by Lemma~\ref{lem:seed-completeness}. For each seed and each
primitive mechanical pattern, compare its word trace with both signed
ray entries $u_{\eps m}$, where
\begin{equation}\label{eq:decision-phase-domain}
 \begin{array}{c|c|c}
 &m=hl+t+q\geqslant1&\text{range of }q\\ \hline
 l\geqslant1&&-\lceil h/2\rceil\leqslant q\leqslant\lfloor3h/2\rfloor\\
 l=0&&-\lceil h/2\rceil\leqslant q\leqslant\lfloor5h/2\rfloor.
 \end{array}
\end{equation}
These are the original phase windows of Section~\ref{sec:gap}.
In particular, the $q$ used here is not changed by the internal
spectral-state interchange in the minimum-product proof.

The domain is a convenient superset. To recover exactly the original
positive-gap family, put $n=\eps m$ and $d=l-n$. If $d\geqslant1$,
retain $(k,t)=(d,t)$. If $d\leqslant-2$, reverse the ray and retain
$(k,t)=(-d-1,h-t)$. Discard $d=0,-1$, which give no original
$k\geqslant1$ representation by this normalization. No equalities are
removed because their trace or seed label has occurred before.

For the positive branch, a retained maximum word cannot have length
$j\geqslant J_++1=2h+4$, since
$$
 (5/2)^{2h+4}>3\cdot4^{h+1}>3H_+.
$$
For the minimum branch, $j\geqslant J_0+1$ would give
$$
 (5/2)^j>2^j\geqslant4(100h^2)^{h+1}>3H_0.
$$
Lemma~\ref{lem:seed-count} gives $S_+,S_0$. There are at most $h-1$
patterns, two signs, and respectively $2h+1$ and $3h+1$ offsets.
This proves \eqref{eq:decision-count}. One may replace $h-1$ by
Euler's totient $\varphi(h)$ when estimating the size of the domain more sharply.

For completeness, seed generation itself has polynomial cost. Each
retained node creates at most two first rejected children, whose labels
are $O(H^2)$. It never explores their descendants. Every entry of a
positive Cohn matrix is at most its trace: the first factor gives row
dominance and the last gives column dominance. Seed matrices and their
inverses therefore have entries $O(H)$. In the positive branch, all
needed exponents are $O(h^2)$ and $\log H_+=O(h)$, so literal products
and recurrence values have $O(h^3)$ bits. In the minimum branch the
corresponding bounds are $O(h)$ and $O(h\log h)$, giving
$O(h^2\log h)$ bits. Exact multiplication adds polynomial overhead
per comparison. Finally
$N_+=O(h^5)$ and $N_0=O(h^4\log^2h)=O(h^5)$.
The promised isometry test is established next.
\end{proof}

\subsection{Recognizing every actual isometry}

\begin{proposition}[The complete homology action]\label{prop:isometry-group}
In the root marking \eqref{eq:root-basis}, the homology image of the full
isometry group of $T$ is exactly
\begin{equation}\label{eq:twelve-actions}
 \mathscr I=\{\pm G^i,\ \pm G^iD_0: i=0,1,2\},\qquad
 G=\begin{pmatrix}0&-1\\1&-1\end{pmatrix},\quad
 D_0=\begin{pmatrix}0&1\\1&0\end{pmatrix}.
\end{equation}
Two primitive simple classes with root homology columns $v,w$ are
isometric if and only if $w=Mv$ for some $M\in\mathscr I$.
\end{proposition}

\begin{proof}
The real matrices
$$
 K_0=\begin{pmatrix}0&1\\-1&1\end{pmatrix},\qquad
 R_0=\begin{pmatrix}1&0\\0&-1\end{pmatrix},\qquad
 \Sigma_0=\begin{pmatrix}0&-1\\1&0\end{pmatrix}
$$
normalize the group generated by $E,F$. Direct multiplication gives
$K_0EK_0^{-1}=F$, $K_0FK_0^{-1}=(EF)^{-1}$;
$R_0ER_0^{-1}=F$, $R_0FR_0^{-1}=E$; and
$\Sigma_0E\Sigma_0^{-1}=E^{-1}$, $\Sigma_0F\Sigma_0^{-1}=F^{-1}$.
The negative determinant of $R_0$ realizes an orientation-reversing
hyperbolic isometry. The resulting homology actions are $G,D_0,-I$,
which generate the twelve distinct matrices in \eqref{eq:twelve-actions}.

To see that the list is complete, use the shortest simple curves.
They are exactly the three classes with normalized trace one, whose
signed primitive columns are
$$
 \mathscr S=\{\pm(1,0),\ \pm(0,1),\ \pm(1,1)\}.
$$
Indeed the initial Farey triangle has these three labels, and every
outward step replaces an entry by a new larger maximum. No later
vertex introduces another label one. An isometry permutes $\mathscr S$.
An integral automorphism preserving $\mathscr S$ has at most six choices
for the image of $(1,0)$. Once that image is fixed, it has at most two
choices for the image of $(0,1)$: after one of the displayed actions
returns the first image to $(1,0)$, the condition that the sum of the
images also belong to $\mathscr S$ leaves only $(0,1)$ and $(-1,-1)$.
Thus there are at most twelve actions, all already realized.
Finally primitive homology determines the simple class up to orientation,
and $-I$ is included. This proves the orbit criterion.
\end{proof}

In a generated seed with columns $\alpha,\beta$, an equality row has
columns
\begin{equation}\label{eq:transported-vectors}
 v=\alpha+n\beta,\qquad w=h\alpha+(hl+t)\beta.
\end{equation}
They are primitive because $\det(\alpha,\beta)=\pm1$ and
$\gcd(h,t)=1$. Their coordinates have size $O(h^3)$: seed columns grow
with their Cohn word length, and the coefficients in
\eqref{eq:transported-vectors} are $O(h^2)$. The twelve comparisons
therefore use $O(\log h)$-bit coordinates. This completes the algorithm
in Theorem~\ref{thm:finite-decision} and distinguishes an equal trace
from an equal trace explained by a surface isometry.
